\chapter{Обзорная часть}

% \subsection{Эллиптически контурированные распределения}

% \textbf{В терминах плотности}

% Случайный вектор $\mathbf{x}\in \mathbb{R}^d$ имеет \textbf{эллиптически контурированное распределение}~\cite{johnson1972distributions}, обозначаемое $\mathrm{EC}_d(\boldsymbol{\mu}, \boldsymbol{\Sigma}, g)$ и также называемое эллиптически симметричным распределением, если $\mathbf{x}$ имеет  плотность вида
% \begin{equation*}
% f(\mathbf{z}) = k_d |\boldsymbol{\Sigma}|^{-1/2} g\left((\mathbf{z} - \boldsymbol{\mu})^\rmT \boldsymbol{\Sigma}^{-1} (\mathbf{z} - \boldsymbol{\mu})\right), 
% \end{equation*}
% где $k_d$ --- нормализующая константа, $\boldsymbol\mu \in \mathbb{R}^d$ --- некоторый параметр положения, $\boldsymbol{\Sigma} \in \mathbb{R}^{d\times d}$ --- симметричная положительно определённая матрица. 

% \vspace{1.5ex}

% \textbf{В терминах характеристических функций}

% Случайный вектор $\mathbf{x}\in \mathbb{R}^d$ имеет \textbf{эллиптически контурированное распределение}~\cite{cambanis1981theory}, если его характеристическая функция $\phi$ удовлетворяет функциональному уравнению для каждого вектора столбца $\mathbf{t}$
% \begin{equation*}
% \phi_{\mathbf{X}}(\mathbf{t}) = \exp(i\mathbf{t}^\rmT \boldsymbol{\mu}) \psi(\mathbf{t}^\rmT \boldsymbol{\Sigma} \mathbf{t}),
% \end{equation*}
% для некоторой скалярной функции $ \psi $. 

% \subsubsection*{Некоторые свойства $\mathrm{EC}$ распределений}

% \begin{enumerate}
%     \item Если существуют вторые моменты, тогда 
%      $$
%      \mathrm{E}(\mathbf{x}) = \boldsymbol{\mu}, \quad \mathrm{Cov}(\mathbf{x}) = c_x \boldsymbol{\Sigma}, \ \text{где} \ c_x = -2 \psi'(0) .
%      $$
%      \item Квадрат расстояния Махаланобиса для $\mathrm{EC}$ имеет плотность~\cite{johnson1987multivariatestatistical}
%      $$
%      h(u) = \frac{\pi^{d/2}}{\Gamma(d/2)} k_d u^{d/2-1}g(u). 
%      $$
%      \item Для $c > 0 \; \mathrm{EC}_d(\boldsymbol{\mu}, c\mathbf{I}, g)$  распределение является сферическим относительно $ \boldsymbol{\mu}$, где $ \mathbf{I} $ --- единичная матрица размера $d \times d$.
%     \end{enumerate}
    
% \textbf{Примеры}  

% \begin{enumerate}
%     \item Многомерное нормальное распределение $N_d(\boldsymbol{\mu}, \boldsymbol{\Sigma})$ является эллиптически контурированным и имеет $k_d = (2\pi)^{-d/2} $, $ \psi(u) = \exp\left(-\frac{u}{2}\right) $ и $ h(u) $ является плотностью распределения хи-квадрат к $d$ степенями свободы $\chi^2_d$.

%     \item Предположим, что  $\mathbf{x}$  является смесью многомерных нормальных с одним и тем же средним и пропорциональными ковариационными матрицами. То есть, пусть
% \begin{equation*}
% \mathbf{x} \sim (1 - \gamma)N_d(\boldsymbol{\mu}, \boldsymbol{\Sigma}) + \gamma N_d(\boldsymbol{\mu}, c\boldsymbol{\Sigma})
% \end{equation*}
% где $ c > 0 $ и $ 0 < \gamma < 1 $. 

%  Смесь также является эллиптически контурированным распределением с $k_d = (2\pi)^{-d/2} $ и $ \psi(u) = (1 - \gamma)\exp\left(-\frac{u}{2}\right) +\gamma \exp\left(-\frac{u}{2c}\right) $.
% \end{enumerate}

\section{Робастное оценивание многомерного параметра положения}

Классические оценки параметра положения, такие как выборочное среднее, обладают высокой чувствительностью к выбросам и отклонениям от предположений о распределении данных. В связи с этим возникает необходимость использования робастных методов оценивания.

В одномерном случае естественной робастной альтернативой среднему является медиана. Однако в многомерном случае отсутствует естественный порядок точек, аналогичный порядку на прямой, поэтому понятие многомерной медианы определяется неоднозначно. В результате было предложено несколько различных конкурирующих определений медианы на многомерный случай. Эти методы отличаются как своими теоретическими свойствами, так и вычислительной сложностью.

Одним из важных свойств многомерных оценок параметра положения и разброса является эквивариантность относительно преобразований данных. Данное свойство означает, что при преобразовании исходной выборки оценка должна изменяться согласованным образом. Поэтому перед рассмотрением конкретных многомерных медиан введем соответствующие определения.

Пусть данные представлены в виде $\mathbf{X} = (\mathbf{x}_1, \ldots, \mathbf{x}_n)^{\mathrm{T}} \in \mathbb{R}^{n \times d}$ --- выборка размера $n$, где каждое наблюдение $\mathbf{x}_i \in \mathbb{R}^d$.

Рассмотрим аффинное преобразование данных:
\begin{equation}\label{eq:affine_transform}
    \mathbf{Z} = \mathbf{X} \mathbf{A}^{\mathrm{T}}+ \mathbf{B},
\end{equation}
где $\mathbf{A}$ --- невырожденная матрица размера $d \times d$, а $\mathbf{B} = \mathbf{1}_n\mathbf{b}^{\mathrm{T}}$, где $\mathbf{1}_n$ --- вектор из $n$ единиц, $\mathbf{b}$ --- вектор размерности $d$.
\medskip

Многомерные оценки положения и разброса $(T, \mathbf{C})$, где $T\in\mathbb{R}^d$, $\mathbf{C}$ --- симметричная положительно определённая матрица, называются \textbf{аффинно эквивариантными}, если справедливы следующие соотношения:
\begin{align}
    T( \mathbf{Z}) &= \mathbf{A}T(\mathbf{X}) + \mathbf{b}, \label{eq:affine_eq_t} \\
    \mathbf{C}( \mathbf{Z}) &=  \mathbf{A} \mathbf{C}(\mathbf{X}) \mathbf{A}^{\mathrm{T}}. \label{eq:affine_eq_c}
\end{align}

% Оценки, которые используют случайные подвыборки или проекции, не являются аффинно эквивариантными, так как они часто меняются при повторном вычислении. Такие оценки иногда могут быть сделаны псевдо-аффинно эквивариантными, если зафиксировать псевдо-случайную последовательность (зафиксировать начальное значение и генератор псевдо-случайных чисел) \cite{olive2017robust}.

Если же для пары из оценки положения и разброса $(T, \mathbf{C})$ соотношения \eqref{eq:affine_eq_t} и \eqref{eq:affine_eq_c} выполняются только для ортогональной матрицы $\mathbf{A}$, то такие оценки называются \textbf{эквивариантными относительно поворота}.

Свойство эквивариантности является важным, поскольку позволяет получать согласованные результаты при изменении системы координат, масштабировании или вращении данных. В частности, при сравнении различных многомерных медиан важно учитывать не только их робастные свойства, но и то, относительно каких преобразований они являются эквивариантными.


\subsection{Геометрическая медиана}

Одним из наиболее известных обобщений одномерной медианы является геометрическая медиана, также называемая пространственной медианой или $L_1$-медианой.

Геометрическая медиана определяется как решение задачи минимизации:
\begin{equation*}
M_s(\mathbf{X})= \arg\min_{\mathbf{x}\in \mathbb{R}^d} \sum\limits_{i=1}^{n}\|\mathbf{x}_i-\mathbf{x}\|.
\end{equation*}

Иными словами, геометрическая медиана является точкой, для которой сумма евклидовых расстояний до наблюдений выборки минимальна. 

Геометрическая медиана является робастной оценкой параметра положения и обладает сравнительно хорошими вычислительными свойствами. Кроме того, она эквивариантна относительно вращений и сдвигов. Благодаря этому она широко используется на практике.

% В предыдущем исследовании было показано, что именно геометрическая медиана является главным конкурентом предложенного алгоритма по точности и скорости работы.


\subsection{Медиана Тьюки}

Другим известным подходом является медиана Тьюки, основанная на понятии глубины точки относительно выборки.

Пусть $H$ --- множество всех замкнутых полупространств в $\mathbb{R}^d$. Для каждого полупространства определяется доля наблюдений выборки, попадающих в него. Глубиной точки называется минимальная доля наблюдений среди всех полупространств, содержащих данную точку. Медианой Тьюки называется множество точек максимальной глубины.

Медиана Тьюки обладает сильными робастными свойствами и является аффинно-эквивариантной. Благодаря этому она считается одной из наиболее теоретически обоснованных многомерных медиан.

Однако главным недостатком данного подхода является высокая вычислительная сложность. С ростом размерности пространства и объема выборки вычислительные затраты быстро возрастают, что ограничивает практическое применение медианы Тьюки, особенно в задачах с большим числом наблюдений или высокой размерностью.


\subsection{Медиана Оджа}

Еще одним обобщением многомерной медианы является медиана Оджа.

Пусть рассматриваются всевозможные симплексы, построенные на точках выборки и некоторой точке $\mathbf{x}$. Тогда медиана Оджа определяется как точка, минимизирующая сумму объемов таких симплексов:
\begin{equation*}
M_o(\mathbf{X})=  \arg\min_{\mathbf{x}\in \mathbb{R}^d} \sum\limits_{i_1 <\ldots<i_d} \mathrm{V}(\mathbf{x}_{i_1},\ldots,\mathbf{x}_{i_d}, \mathbf{x}),
\end{equation*}
где $\mathbf{x}_{i_1},\ldots,\mathbf{x}_{i_d},\mathbf{x}$ --- вершины, которыми определяется симплекс, $\{1\leqslant i_1 <\ldots<i_d\leqslant n\}$ --- набор из всех упорядоченных $d$-кортежей из $\{1,\ldots,n\},\; 1\leqslant d \leqslant n$.

Как и медиана Тьюки, медиана Оджа является аффинно-эквивариантной. Это делает ее привлекательной с теоретической точки зрения. Однако вычислительная сложность данного метода также очень высока. Количество рассматриваемых симплексов быстро растет с увеличением размерности пространства и объема выборки, что делает практическое вычисление медианы Оджа трудоемким.


\subsection{Проекционная медиана}

Проекционная медиана, введенная в \cite{durocher_projection_2009}, представляет собой еще один подход к построению робастной оценки многомерного параметра положения. 
\begin{equation}\label{eq:projmed}
M_p(\mathbf{X}) = d\frac{\int_{\{\mathbf{a}: \, \mathbf{a}^{\mathrm{T}} \mathbf{a} = 1\}} \text{med}(\mathbf{X}_{\mathbf{a}}) \textrm{d}\mathbf{a}}{\int_{\{\mathbf{a}: \, \mathbf{a}^{\mathrm{T}} \mathbf{a} = 1\}} \textrm{d}\mathbf{a}},
\end{equation}
где $\mathbf{a} \in \mathbb{R}^{d}$, $\text{med}(\mathbf{X}_{\mathbf{a}})$ обозначает (одномерную) медиану множества точек $\mathbf{X}$, спроектированных на направление вдоль векторa $\mathbf{a}$, проходящее через начало координат. Эта оценка является эквивариантной относительно поворотов, однако вычислить её аналитически из-за наличия интеграла по единичной сфере в общем случае невозможно.

В \cite{durocher_projection_2017} было получено приближение к \eqref{eq:projmed} по методу Монте-Карло:
\begin{equation*}
\widetilde M_{p, K}(\mathbf{X}) = \frac{d}{K} \sum_{i=1}^{K} \text{med}(\mathbf{X}_{\boldsymbol{ \omega}_i}),
\end{equation*}
где $ \boldsymbol{\omega}_1, \ldots, \boldsymbol{\omega}_K$ --- независимые одинаково распределенные случайные величины, а $\Omega_d$ --- равномерное распределение на единичной сфере $\mathbf{a}^{\mathrm{T}} \mathbf{a} = 1$.

\subsubsection{Метод Yamm}

Введем обозначения: пусть $\boldsymbol{\mu}  \in \mathbb{R}^d $ --- вектор сдвига,  $\mathbf{y} = \mathbf{y}(\boldsymbol{\mu}, \mathbf{a})$ --- проекция $\mathbf{X}$ на $\mathbf{a}$ после того, как начало координат переместили в $\boldsymbol{\mu}$: $\mathbf{y} = (\mathbf{X} - 1_n \boldsymbol{\mu}^{\mathrm{T}} ) \mathbf{a}$, где $\mathbf{y} \in \mathbb{R}^n$, и $m_{\mathbf{X}}(\boldsymbol{\mu}, \mathbf{a}) = m(\mathbf{y})$ --- одномерная медиана спроецированных точек~$\mathbf{y}$.


Определим интеграл
\begin{equation}\label{eq:yamm_tf}
M_{\mathbf{X}}(\boldsymbol{\mu}) = \int_{\{\mathbf{a}: \, \mathbf{a}^{\mathrm{T}} \mathbf{a} = 1\}} m_{\mathbf{X}}(\boldsymbol{\mu}, \mathbf{a})^2 \textrm{d} \,\mathbf{a},
\end{equation}
и поставим задачу минимизации:
\begin{equation}\label{eq:yamm}
\hat{\boldsymbol{\mu}} = \arg\min_{\boldsymbol{\mu} \in \mathbb{R}^{d}} M_{\mathbf{X}}(\boldsymbol{\mu}).
\end{equation}

Оценка \eqref{eq:yamm} называется оценкой \textbf{yamm} \cite{chen_new_2020}. Для неё справедлив следующий результат.

\newpage

\begin{theorem}[\cite{chen_new_2020}]\label{th:equiv}
Для любой конечной выборки $\mathbf{X}\in \mathbb{R}^{d}$ оценка yamm \eqref{eq:yamm} эквивалентна проекционной медиане \eqref{eq:projmed}.
\end{theorem}

Из Теоремы \ref{th:equiv} также следует, что оценка yamm, как эквивалентная, и проекционная медиана, являются эквивариантными относительно поворотов. Известно, однако, что проекционная медиана не является аффинно-эквивариантной.
 
Для оценивания по методу yamm вместо оптимизации целевой функции \eqref{eq:yamm_tf}, которую, как и функцию \eqref{eq:projmed}, трудно вычислить аналитически, делают следующее: разыгрывают $\boldsymbol{\omega}_1$, \ldots, $\boldsymbol{\omega}_K$ --- независимые одинаково распределенные из $\Omega_d$, после чего следующий функционал оптимизируют по $\boldsymbol{\mu}$, используя оптимизатор BFGS \cite{nocedal2006numerical}:
\begin{equation*}
\tilde{\boldsymbol{\mu}} = \arg\min_{\boldsymbol{\mu \in \mathbb{R}^{d}}} \sum_{i=1}^K m_{\mathbf{X}}(\boldsymbol{\mu}, \mathbf{\boldsymbol{ \omega}}_i)^2.
\end{equation*}

Таким образом, вместо оптимизации исходного функционала используется оптимизация его приближения по методу Монте-Карло из-за трудности аналитического вычисления \eqref{eq:yamm_tf}.

Интерес к данной оценке связан с тем, что она позволяет использовать одномерные медианы на случайных проекциях для построения робастной многомерной оценки положения. Кроме того, проекционная медиана во многом рассматривается как попытка объединить хорошие робастные свойства многомерных медиан с более приемлемой вычислительной сложностью.


\section{Робастные оценки положения и разброса на основе алгоритмов концентрации}\label{sec:concentration}

В предыдущем разделе рассматривались робастные оценки параметра положения. Однако во многих задачах многомерного анализа требуется оценивать не только центр распределения, но и матрицу разброса. Поэтому далее кратко рассмотрим методы, ориентированные на совместное робастное оценивание параметра положения и матрицы разброса.

Классическими оценками положения и разброса являются выборочное среднее и выборочная ковариационная матрица. Однако обе эти оценки чувствительны к выбросам, поэтому при их наличии  могут давать существенно искаженные результаты.

Для повышения робастности используются различные приближенные алгоритмы, среди которых важное место занимают алгоритмы концентрации. Их основная идея заключается в итеративном уточнении оценок. Процесс начинается с некоторой \textit{начальной точки}, представляющей предварительную оценку параметров распределения. Далее на каждом шаге формируется новая оценка, называемая \textit{аттрактором}, которая вычисляется по подмножеству наблюдений, наиболее согласованных с текущей оценкой.

Например, пусть в качестве начальной точки используется классическая оценка $(\bar{\mathbf{x}}, \widehat{\mathbf{S}})$ --- пара из выборочного среднего и выборочной ковариационной матрицы. Тогда аттрактором может быть классическая оценка $(T_1, \mathbf{C}_1)$, вычисленная по той половине выборки, для которой расстояния Махаланобиса оказываются наименьшими. Этот алгоритм использует один шаг концентрации, но этот процесс может быть продлён и до $k$ шагов концентрации, последовательно образуя оценку $(T_k, \mathbf{C}_k)$.

Если используется более одного аттрактора, то необходимо выбрать некоторый критерий для получения финальной оценки. Если каждый аттрактор $(T_{k,j}, \mathbf{C}_{k,j})$ является классической оценкой, применяемой примерно к $n/2$ случаев, то обычно используется критерий минимального определителя ковариационной матрицы (MCD): выбирается аттрактор, имеющий наименьшее значение $\det(\mathbf{C}_{k,j})$, где $j=1, \ldots, k$.

\vspace{1.5ex}

\textbf{MCD}

Оценка MCD основывается на подмножестве $J_0$, содержащем $c_n \approx n/2$ наблюдений. Это подмножество выбирается таким образом, чтобы его выборочная ковариационная матрица имела наименьший определитель среди всех возможных $C_n^{c_n}$ подмножеств размера $c_n$ (здесь $C_n^i = \frac{n!}{i! (n-i)!}$ ---  биномиальный коэффициент). Выборочное среднее и выборочная ковариационная матрица, рассчитанные по наблюдениям из $J_0$, обозначаются $T_{MCD}$ и $\mathbf{C}_{MCD}$ соответственно. В результате, пара $(T_{MCD}, \mathbf{C}_{MCD})$ формирует оценку \textit{минимального определителя ковариационной матрицы}.

Точный перебор всех возможных подмножеств $C_n^{c_n}$ требует слишком больших вычислительных затрат даже для умеренных объемов выборки  $n$. Поэтому на практике применяются различные приближенные алгоритмы.

\vspace{1.5ex}

\textbf{Алгоритм концентрациии}

Алгоритмы концентрации позволяют приближенно вычислять робастные оценки, используя ограниченное число начальных точек и итеративное уточнение.

Алгоритм концентрации начинается с элементарного набора данных $J$ размером $d+1$, где $d$ --- размерность пространства. Элементарной начальной точкой является выборочное среднее и выборочная ковариационная матрица, рассчитанные по наблюдениям из $J$.

Обозначим $j$-ую начальную точку как $(T_{-1,j}, \mathbf{C}_{-1,j})$. Для каждой такой начальной точки вычисляются расстояния Махаланобиса для всех $n$ наблюдений выборки. На следующей итерации формируется оценка $(T_{0,j}, \mathbf{C}_{0,j})=(\bar{\mathbf{x}}_{0,j}, \widehat{\mathbf{S}}_{0,j})$, которая рассчитывается на основе $c_n \approx n/2$ наблюдений с наименьшими расстояниями. Этот итеративный процесс может выполняться до $k$ шагов концентрации, в результате чего получается последовательность оценок $(T_{-1,j}, \mathbf{C}_{-1,j}),$ $(T_{0,j}, \mathbf{C}_{0,j}),\ldots, (T_{k,j}, \mathbf{C}_{k,j})$. 

Итоговая оценка $(T_{k,j}, \mathbf{C}_{k,j})$ представляет собой аттрактор. Если используется $K_n$ начальных точек, то финальный аттрактор $(T_A, \mathbf{C}_A)$ выбирается согласно некоторому критерию. Часто применяется критерий минимального определителя: аттрактор с наименьшим $\det(\mathbf{C}_{k,j})$ выбирается в качестве итоговой оценки.

Таким образом, методы концентрации фактически реализуют итеративное уточнение области, которая считается <<наиболее согласованной>> частью выборки. Благодаря этому удается уменьшить влияние выбросов и повысить устойчивость оценки.

\newpage

\textbf{Fast-MCD, FCH, MBA, RFCH, RMVN}\footnote{В данной работе методы, основанные на алгоритмах концентрации, не рассматриваются как прямые конкуренты предложенного алгоритма при сравнении оценок параметра положения. Однако они важны для дальнейшего изложения, поскольку демонстрируют общий подход к совместному робастному оцениванию положения и разброса. Именно эта идея будет использована далее при построении оценки матрицы разброса.}

Кратко опишем робастные оценки, применяемые на практике, основанные  на алгоритмах концентрации.

\begin{itemize}
    \item Оценка \textbf{Fast-MCD}~\cite{hubert2008high}  использует классическую оценку, примененную к $K = 500$ случайно выбранным элементарным наборам из $d+1$ точек в качестве стартовых. Финальная оценка выбирается по критерию минимального определителя.

    \item  Оценки \textbf{FCH} и \textbf{MBA} используют два аттрактора. Первый аттрактор получается с использованием оценки DGK~\cite{devlin1981robust}, где в качестве начальной точки применяется классическая оценка $(\bar{\mathbf{x}}, \widehat{\mathbf{S}})$. Второй аттрактор формируется с помощью оценки медианного шара MB~\cite{olive2004resistant}, начальной точкой для которой служит $(\text{MED}(\textbf{W}), \textbf{I}_d)$, где $\text{MED}(\textbf{W})$ --- покоординатная медиана, а $\textbf{I}_d$ --- единичная матрица размера $d \times d$.

    \item  Оценка \textbf{RFCH} использует два шага перевзвешивания для повышения эффективности~\cite{olive2010robust}, а оценка \textbf{RMVN} применяет модифицированный метод перевзвешивания~\cite{olive2010robust}.
\end{itemize} 
